\documentclass[UTF8,a4paper,11pt]{ctexart}
\usepackage[margin=22mm,headheight=14pt,footskip=12mm]{geometry}
\usepackage{amsmath,amssymb,array,longtable,booktabs,tikz,needspace,enumitem,fancyhdr,hyperref}
\usepackage{microtype}
\hypersetup{colorlinks=false,pdfborder={0 0 0},pdfauthor={Good Learning},bookmarks=false}
\newfontfamily\circledfont{Segoe UI Symbol}
\newcolumntype{P}[1]{>{\raggedright\arraybackslash}p{#1}}
\setlength{\LTleft}{0pt}\setlength{\LTright}{0pt}
\setlength{\parindent}{0pt}\setlength{\parskip}{.65em}
\linespread{1.18}
\raggedbottom
\setcounter{secnumdepth}{0}
\setlist{leftmargin=2.2em,itemsep=.4em,topsep=.5em}
\ctexset{section={format=\Large\bfseries\sffamily,beforeskip=1.35em,afterskip=.65em},
subsection={format=\large\bfseries\sffamily,beforeskip=1em,afterskip=.45em}}
\newcommand{\chaptertitle}[1]{\Needspace{8\baselineskip}\noindent{\LARGE\bfseries\sffamily #1}\par\vspace{4mm}\hrule\vspace{6mm}}
\newcommand{\CompactQuestion}[2]{\par\noindent\textbf{#1.}\ #2\par\vspace{.8em}}
\newcommand{\PracticeQuestion}[3]{\par\Needspace{4\baselineskip}\noindent\begin{minipage}{\linewidth}\textbf{#1.}\ #2\par\noindent\rule{0pt}{#3mm}\par\end{minipage}\par}
\newcommand{\CompactQuestionColon}[2]{\par\noindent\textbf{#1：}\ #2\par\vspace{.8em}}
\newcommand{\PracticeQuestionColon}[3]{\par\Needspace{4\baselineskip}\noindent\begin{minipage}{\linewidth}\textbf{#1：}\ #2\par\noindent\rule{0pt}{#3mm}\par\end{minipage}\par}
\pagestyle{fancy}\fancyhf{}\fancyfoot[C]{\thepage}\renewcommand{\headrulewidth}{0pt}
\begin{document}
\thispagestyle{empty}
\noindent{\sffamily\footnotesize 知识结构}\hfill{\sffamily\footnotesize 学习导图}\par
\vspace{7mm}
\noindent{\sffamily\fontsize{25}{30}\selectfont\bfseries 应用随机过程 · 知识结构}\par
\vspace{4mm}\noindent\rule{\linewidth}{.85pt}\par\vspace{6mm}
\begin{center}
\begin{tikzpicture}[x=1cm,y=1cm,>=stealth,every node/.style={font=\sffamily}]
\path[use as bounding box] (-8.1,-18.7) rectangle (8.1,.4);
\node[font=\sffamily\Large\bfseries] at (0,-.65) {应用随机过程 · 七章结构};
\node[font=\sffamily\small,text=black!55] at (0,-1.35) {按章节编号学习并自检};
\draw[line width=.85pt] (0,-1.9) -- (0,-2.45);
\draw[line width=.85pt] (-4,-2.45) -- (4,-2.45);
\draw[line width=.85pt] (-4,-2.45) -- (-4,-2.8);
\draw[line width=.85pt] (4,-2.45) -- (4,-2.8);
\node[font=\sffamily\bfseries] at (-4,-3.15) {基础与计数过程};
\node[font=\sffamily\bfseries] at (4,-3.15) {马氏链与 Brown 运动};
\node[anchor=north west,font=\sffamily\fontsize{22}{24}\selectfont\bfseries,text=black!50] at (-7.4,-4.3) {01};
\node[anchor=north west,text width=5.9cm,align=left,font=\sffamily\large\bfseries] at (-6.25,-4.28) {第1章 预备知识};
\node[anchor=north west,text width=5.9cm,align=left,font=\sffamily\small] at (-6.25,-5.05) {1.1—1.5 · 001—031\\[.25em]特征函数与条件期望是核心};
\node[anchor=west,font=\sffamily\footnotesize,text=black!60] at (-6.25,-6.35) {最高优先 · 5 小时};
\draw[black!35,line width=.4pt] (-7.4,-6.65) -- (-0.25,-6.65);
\node[anchor=north west,font=\sffamily\fontsize{22}{24}\selectfont\bfseries,text=black!50] at (-7.4,-8.399999999999999) {02};
\node[anchor=north west,text width=5.9cm,align=left,font=\sffamily\large\bfseries] at (-6.25,-8.379999999999999) {第2章 基本概念};
\node[anchor=north west,text width=5.9cm,align=left,font=\sffamily\small] at (-6.25,-9.149999999999999) {2.1—2.3 · 032—042\\[.25em]有限维分布 · 平稳性 · 增量};
\node[anchor=west,font=\sffamily\footnotesize,text=black!60] at (-6.25,-10.45) {最高优先 · 2 小时};
\draw[black!35,line width=.4pt] (-7.4,-10.749999999999998) -- (-0.25,-10.749999999999998);
\node[anchor=north west,font=\sffamily\fontsize{22}{24}\selectfont\bfseries,text=black!50] at (-7.4,-12.5) {03};
\node[anchor=north west,text width=5.9cm,align=left,font=\sffamily\large\bfseries] at (-6.25,-12.48) {第3章 Poisson 过程};
\node[anchor=north west,text width=5.9cm,align=left,font=\sffamily\small] at (-6.25,-13.25) {3.1—3.3 · 043—060\\[.25em]条件分布与二阶矩为高频};
\node[anchor=west,font=\sffamily\footnotesize,text=black!60] at (-6.25,-14.55) {最高优先 · 5 小时};
\draw[black!35,line width=.4pt] (-7.4,-14.85) -- (-0.25,-14.85);
\node[anchor=north west,font=\sffamily\fontsize{22}{24}\selectfont\bfseries,text=black!50] at (0.8,-4.3) {04};
\node[anchor=north west,text width=5.9cm,align=left,font=\sffamily\large\bfseries] at (1.95,-4.28) {第4章 更新过程};
\node[anchor=north west,text width=5.9cm,align=left,font=\sffamily\small] at (1.95,-5.05) {4.1—4.4 · 062—081\\[.25em]更新方程与更新定理};
\node[anchor=west,font=\sffamily\footnotesize,text=black!60] at (1.95,-6.35) {中高优先 · 4 小时};
\draw[black!35,line width=.4pt] (0.8,-6.65) -- (7.95,-6.65);
\node[anchor=north west,font=\sffamily\fontsize{22}{24}\selectfont\bfseries,text=black!50] at (0.8,-8.399999999999999) {05};
\node[anchor=north west,text width=5.9cm,align=left,font=\sffamily\large\bfseries] at (1.95,-8.379999999999999) {第5章 Markov 链};
\node[anchor=north west,text width=5.9cm,align=left,font=\sffamily\small] at (1.95,-9.149999999999999) {5.1—5.5 · 082—120\\[.25em]平稳分布与吸收概率};
\node[anchor=west,font=\sffamily\footnotesize,text=black!60] at (1.95,-10.45) {最高优先 · 8 小时};
\draw[black!35,line width=.4pt] (0.8,-10.749999999999998) -- (7.95,-10.749999999999998);
\node[anchor=north west,font=\sffamily\fontsize{22}{24}\selectfont\bfseries,text=black!50] at (0.8,-12.5) {06};
\node[anchor=north west,text width=5.9cm,align=left,font=\sffamily\large\bfseries] at (1.95,-12.48) {第6—7章 鞅与 Brown};
\node[anchor=north west,text width=5.9cm,align=left,font=\sffamily\small] at (1.95,-13.25) {121—160\\[.25em]鞅的停时 · 反射原理};
\node[anchor=west,font=\sffamily\footnotesize,text=black!60] at (1.95,-14.55) {高优先 · 7 小时};
\draw[black!35,line width=.4pt] (0.8,-14.85) -- (7.95,-14.85);
\node[font=\sffamily\small,text=black!55] at (0,-17.95) {编号对应知识清单 · 页码见正文};
\end{tikzpicture}
\end{center}
\clearpage

\chaptertitle{应用随机过程 · 知识清单（第二版）}
教材：张波、商豪、邓军《应用随机过程》（第6版），中国人民大学出版社，2023 年 10 月。考试：2026-11-15。范围：教材第 1—7 章（第 001—160 页）。配套材料：李英老师授课笔记（21 页，第 1—4 页对应第 1—2 章）、2023 级随机过程必修期末考试试题及手写解答（2 页）。

本版已把条目落到具体章节与页码，不再使用通用范围。证据优先级：教材正文 = 真题 > 课堂笔记 > 往年经验。真题只见过 2023 级一份，题型权重仍标“样本有限”，不外推为精确分值。

\Needspace{6\baselineskip}
\section{一、范围与证据}
{\small\renewcommand{\arraystretch}{1.22}\setlength{\tabcolsep}{3.5pt}
\begin{longtable}{P{0.17\linewidth}P{0.48\linewidth}P{0.28\linewidth}}
\toprule
\bfseries 项目 & \bfseries 依据 & \bfseries 说明 \\
\midrule
\endfirsthead
\toprule
\bfseries 项目 & \bfseries 依据 & \bfseries 说明 \\
\midrule
\endhead
章节范围 & 教材目录 & 第 1 章 001、第 2 章 032、第 3 章 043、第 4 章 062、第 5 章 082、第 6 章 121、第 7 章 144 \\
\addlinespace[.15em]
真题题型 & 2023 级期末卷 & 全卷以大题为主：特征函数与和的分布、Poisson 条件概率、Poisson 二阶矩、分布/最大值、Markov 链平稳分布、鞅、期望方程、写出 Poisson 与 Brown 定义 \\
\addlinespace[.15em]
教师侧重 & 李英笔记 & 大量篇幅在特征函数、条件期望、Poisson 过程条件分布、Markov 链平稳分布、鞅的停时，与真题方向一致 \\
\addlinespace[.15em]
未覆盖 & 第 8 章随机积分、9—11 章 & 不在范围内，不学、不做题 \\
\addlinespace[.15em]
\bottomrule
\end{longtable}
}
\Needspace{6\baselineskip}
\section{二、第 1 章 预备知识（001—031）}
{\small\renewcommand{\arraystretch}{1.22}\setlength{\tabcolsep}{3.5pt}
\begin{longtable}{P{0.08\linewidth}P{0.40\linewidth}P{0.29\linewidth}P{0.16\linewidth}}
\toprule
\bfseries 编号 & \bfseries 知识点 & \bfseries 教材位置与依据 & \bfseries 优先级/状态 \\
\midrule
\endfirsthead
\toprule
\bfseries 编号 & \bfseries 知识点 & \bfseries 教材位置与依据 & \bfseries 优先级/状态 \\
\midrule
\endhead
1.1-a & 概率空间 $(\Omega,\mathcal{F},P)$、事件域、概率的可列可加性 & 1.1（001—007）；笔记第 1 页 & 中/待测 \\
\addlinespace[.15em]
1.2-a & 随机变量与分布函数：三性质、区间概率、单点概率、Riemann–Stieltjes 积分 & 1.2（007—013）；笔记第 1 页例 1—2 & 高/待测 \\
\addlinespace[.15em]
1.3-a & 期望、方差、协方差；矩母函数 $\psi(t)=E(e^{tX})$ 与求矩公式 & 1.3（013—021）；笔记第 1 页 & 高/待测 \\
\addlinespace[.15em]
1.3-b & 特征函数 $\varphi(t)=E(e^{itX})$：性质、独立和乘积、唯一性定理、五个常见分布 & 1.3；2023 卷第一道大题考过 & 最高/待测 \\
\addlinespace[.15em]
1.3-c & 由 $\varphi^{(k)}(0)=i^kE(X^k)$ 求高阶矩 & 1.3；笔记第 2 页题 3 & 高/待测 \\
\addlinespace[.15em]
1.4-a & 四种收敛及其推出关系 & 1.4（021—023）；低频 & 低/待测 \\
\addlinespace[.15em]
1.5-a & 条件期望定义、全期望公式、独立性、条件方差、Wald 公式 & 1.5（023—031）；笔记第 2 页题 5 & 最高/待测 \\
\addlinespace[.15em]
\bottomrule
\end{longtable}
}
\Needspace{6\baselineskip}
\section{三、第 2 章 随机过程的基本概念和基本类型（032—042）}
{\small\renewcommand{\arraystretch}{1.22}\setlength{\tabcolsep}{3.5pt}
\begin{longtable}{P{0.08\linewidth}P{0.40\linewidth}P{0.29\linewidth}P{0.16\linewidth}}
\toprule
\bfseries 编号 & \bfseries 知识点 & \bfseries 教材位置与依据 & \bfseries 优先级/状态 \\
\midrule
\endfirsthead
\toprule
\bfseries 编号 & \bfseries 知识点 & \bfseries 教材位置与依据 & \bfseries 优先级/状态 \\
\midrule
\endhead
2.1-a & 随机过程定义、指标集、状态空间、样本函数 & 2.1；笔记第 3 页 & 高/待测 \\
\addlinespace[.15em]
2.1-b & 一维分布 $F_t(x)$ 与其局限 & 2.1；笔记第 3 页 & 中/待测 \\
\addlinespace[.15em]
2.2-a & 有限维分布族、对称性、相容性、Kolmogorov 定理 & 2.2；笔记第 3 页 & 高/待测 \\
\addlinespace[.15em]
2.2-b & 用独立增量求联合分布 & 2.2 例题；笔记第 4 页例 7 & 高/待测 \\
\addlinespace[.15em]
2.3-a & 均值函数、方差函数、自相关函数、协方差函数 & 2.3；笔记第 3—4 页例 5 & 最高/待测 \\
\addlinespace[.15em]
2.3-b & 严平稳与宽平稳、两者关系、正态过程等价 & 2.3；笔记第 3 页 & 最高/待测 \\
\addlinespace[.15em]
2.3-c & 独立增量、平稳增量、独立平稳增量过程 & 2.3；笔记第 4 页例 3 & 最高/待测 \\
\addlinespace[.15em]
2.3-d & 典型过程：部分和、随机电报信号、Poisson 与 Brown 的预告 & 2.3；笔记第 4 页 & 高/待测 \\
\addlinespace[.15em]
\bottomrule
\end{longtable}
}
\Needspace{6\baselineskip}
\section{四、第 3 章 Poisson 过程（043—060）}
{\small\renewcommand{\arraystretch}{1.22}\setlength{\tabcolsep}{3.5pt}
\begin{longtable}{P{0.08\linewidth}P{0.40\linewidth}P{0.29\linewidth}P{0.16\linewidth}}
\toprule
\bfseries 编号 & \bfseries 知识点 & \bfseries 教材位置与依据 & \bfseries 优先级/状态 \\
\midrule
\endfirsthead
\toprule
\bfseries 编号 & \bfseries 知识点 & \bfseries 教材位置与依据 & \bfseries 优先级/状态 \\
\midrule
\endhead
3.1-a & 计数过程、Poisson 过程定义（独立平稳增量 + 强度）、等价刻画 & 3.1（043—049）；笔记第 5 页；真题要求写定义 & 最高/待测 \\
\addlinespace[.15em]
3.1-b & 到达间隔独立同指数、$T_n\sim\Gamma(n,\lambda)$ & 3.1；笔记第 5 页 & 最高/待测 \\
\addlinespace[.15em]
3.1-c & 条件分布：给定 $N(t)=n$ 时到达时刻为均匀次序统计量 & 3.2（049—054）；笔记第 6 页题 3；真题有同类题 & 最高/待测 \\
\addlinespace[.15em]
3.2-a & 到达时刻条件期望 $E[\sum (t-T_i)\mid N(t)=n]$ & 3.2；笔记第 6 页例题 & 高/待测 \\
\addlinespace[.15em]
3.2-b & 二阶矩与协方差 $E[N(s)N(t)]$、$R_N(s,t)$ & 3.2；真题考过 $E[N(s)N(t)]$ & 最高/待测 \\
\addlinespace[.15em]
3.3-a & 非齐次 Poisson 过程与均值函数 $m(t)$ & 3.3（054—060）；笔记第 5 页 & 中/待测 \\
\addlinespace[.15em]
3.3-b & 复合 Poisson 过程与条件期望 & 3.3；笔记第 6 页题 1 & 高/待测 \\
\addlinespace[.15em]
3.3-c & 叠加、稀疏（分解）与次序统计量应用 & 3.3；教材例题 & 中/待测 \\
\addlinespace[.15em]
\bottomrule
\end{longtable}
}
\Needspace{6\baselineskip}
\section{五、第 4 章 更新过程（062—081）}
{\small\renewcommand{\arraystretch}{1.22}\setlength{\tabcolsep}{3.5pt}
\begin{longtable}{P{0.08\linewidth}P{0.40\linewidth}P{0.29\linewidth}P{0.16\linewidth}}
\toprule
\bfseries 编号 & \bfseries 知识点 & \bfseries 教材位置与依据 & \bfseries 优先级/状态 \\
\midrule
\endfirsthead
\toprule
\bfseries 编号 & \bfseries 知识点 & \bfseries 教材位置与依据 & \bfseries 优先级/状态 \\
\midrule
\endhead
4.1-a & 更新过程定义、更新函数 $m(t)$、到达时刻分布 & 4.1（062—065） & 高/待测 \\
\addlinespace[.15em]
4.2-a & 更新方程及其求解（Laplace 变换、逐次代入） & 4.2（065—070） & 高/待测 \\
\addlinespace[.15em]
4.3-a & 更新定理、更新报酬定理、长期平均率 & 4.3（070—078） & 高/待测 \\
\addlinespace[.15em]
4.4-a & 剩余寿命与年龄、交错更新与可用度 & 4.4（078—081） & 中/待测 \\
\addlinespace[.15em]
\bottomrule
\end{longtable}
}
\Needspace{6\baselineskip}
\section{六、第 5 章 Markov 链（082—120）}
{\small\renewcommand{\arraystretch}{1.22}\setlength{\tabcolsep}{3.5pt}
\begin{longtable}{P{0.08\linewidth}P{0.40\linewidth}P{0.29\linewidth}P{0.16\linewidth}}
\toprule
\bfseries 编号 & \bfseries 知识点 & \bfseries 教材位置与依据 & \bfseries 优先级/状态 \\
\midrule
\endfirsthead
\toprule
\bfseries 编号 & \bfseries 知识点 & \bfseries 教材位置与依据 & \bfseries 优先级/状态 \\
\midrule
\endhead
5.1-a & 定义、转移矩阵、C-K 方程、$n$ 步转移 & 5.1（082—091）；真题摸球题第一步 & 最高/待测 \\
\addlinespace[.15em]
5.2-a & 状态分类：常返/瞬态、周期、不可约、闭集 & 5.2（091—097） & 最高/待测 \\
\addlinespace[.15em]
5.3-a & 平稳分布 $\pi P=\pi$、细致平衡、极限定理 & 5.3（097—105）；真题摸球题 & 最高/待测 \\
\addlinespace[.15em]
5.4-a & 首达时间与吸收概率、期望方程（赌徒破产类） & 5.4（105—110）；真题矿工题同类 & 最高/待测 \\
\addlinespace[.15em]
5.4-b & 应用建模：生灭链、M/G/1 嵌入链 & 5.4 & 中/待测 \\
\addlinespace[.15em]
5.5-a & 连续时间 Markov 链、Q 矩阵、Kolmogorov 方程 & 5.5（110—118） & 高/待测 \\
\addlinespace[.15em]
5.5-b & 生灭过程、M/M/1 平稳分布 & 5.5 & 高/待测 \\
\addlinespace[.15em]
\bottomrule
\end{longtable}
}
\Needspace{6\baselineskip}
\section{七、第 6 章 鞅（121—142）}
{\small\renewcommand{\arraystretch}{1.22}\setlength{\tabcolsep}{3.5pt}
\begin{longtable}{P{0.08\linewidth}P{0.40\linewidth}P{0.29\linewidth}P{0.16\linewidth}}
\toprule
\bfseries 编号 & \bfseries 知识点 & \bfseries 教材位置与依据 & \bfseries 优先级/状态 \\
\midrule
\endfirsthead
\toprule
\bfseries 编号 & \bfseries 知识点 & \bfseries 教材位置与依据 & \bfseries 优先级/状态 \\
\midrule
\endhead
6.1-a & 鞅的定义、下鞅/上鞅、条件期望形式 & 6.1（121—126）；真题硬币鞅题 & 最高/待测 \\
\addlinespace[.15em]
6.2-a & 停时、可选停时定理（结论与应用条件） & 6.2（126—136）；真题硬币鞅题 & 最高/待测 \\
\addlinespace[.15em]
6.3-a & 一致可积性 & 6.3（136—137） & 低/待测 \\
\addlinespace[.15em]
6.4-a & 鞅收敛定理 & 6.4（137—140） & 低/待测 \\
\addlinespace[.15em]
6.5-a & 连续鞅、Brown 运动的鞅性质 & 6.5（140—142）、7.3 & 中/待测 \\
\addlinespace[.15em]
\bottomrule
\end{longtable}
}
\Needspace{6\baselineskip}
\section{八、第 7 章 Brown 运动（144—160）}
{\small\renewcommand{\arraystretch}{1.22}\setlength{\tabcolsep}{3.5pt}
\begin{longtable}{P{0.08\linewidth}P{0.40\linewidth}P{0.29\linewidth}P{0.16\linewidth}}
\toprule
\bfseries 编号 & \bfseries 知识点 & \bfseries 教材位置与依据 & \bfseries 优先级/状态 \\
\midrule
\endfirsthead
\toprule
\bfseries 编号 & \bfseries 知识点 & \bfseries 教材位置与依据 & \bfseries 优先级/状态 \\
\midrule
\endhead
7.1-a & 定义与性质：独立增量、正态增量、连续轨道、二次变差 & 7.1（144—147）；真题要求写定义 & 最高/待测 \\
\addlinespace[.15em]
7.2-a & Gauss 过程、协方差结构 & 7.2（147—149） & 中/待测 \\
\addlinespace[.15em]
7.3-a & Brown 运动的鞅性质 & 7.3（149—151） & 高/待测 \\
\addlinespace[.15em]
7.4-a & Markov 性 & 7.4（151—152） & 中/待测 \\
\addlinespace[.15em]
7.5-a & 最大值变量、反射原理、反正弦律 & 7.5（152—156） & 高/待测 \\
\addlinespace[.15em]
7.6-a & Brown 运动的几种变化（漂移、几何 Brown） & 7.6（156—160） & 中/待测 \\
\addlinespace[.15em]
7.7-a & 高维 Brown 运动 & 7.7（160） & 低/待测 \\
\addlinespace[.15em]
\bottomrule
\end{longtable}
}
\Needspace{6\baselineskip}
\section{九、状态与下一步}
\begin{itemize}
\item \textbf{当前进度}：第 1、2 章已出课程（D001、D002，2026-10-25），全部条目为“待测”，等你提交 A1—A8 与 B1—B8 的答案后再写入真实掌握状态。
\item \textbf{最高优先级集中区}：1.3-b（特征函数）、1.5-a（条件期望）、2.3-a/b/c、3.1、3.2、5.1—5.4、6.1—6.2、7.1。这些对应 2023 级真题的绝大部分分值。
\item \textbf{可延后}：1.4（收敛性）、6.3—6.4（一致可积、鞅收敛）、7.7（高维 Brown），只在时间富余时看结论。
\item \textbf{下一步}：完成今天两课并交答案；之后按 `备考计划.md` 进入第 3 章 Poisson 过程。
\end{itemize}
本文件与 `备考计划.md` 互相引用：清单负责“学什么、学到什么程度”，计划负责“何时学、如何验收”。

\end{document}
